\documentclass[12pt]{article}

\usepackage[a4paper, total={6in, 8in}]{geometry}
\usepackage{textcomp}

\begin{document}
\setlength{\parindent}{0pt}

\def \t {\textrightarrow}
\def \v {\vspace{1em}}
\def \bi {\begin{itemize}}
\def \ei {\end{itemize}}
\def \s[#1] {\section*{#1}}
\def \ss[#1] {\subsection*{#1}}
\def \sss[#1] {\subsubsection*{#1}}

\s[Recuperare 24 febbraio]
\s[Italia dopo 43]
Nel 43 viene firmato l'armistizio \t il paese è diviso in due \\
Inizia il fenomeno della resistenza, da cui usciranno tutti i grandi partiti del dopoguerra \t sara un esperienza di difesa, ma anche politica \\
La resistenza italiana è complessa perche è una guerra di patrioti per liberare il paese dai tedeschi, ma è anche una guerra civile (fascisti contro partigiani, i fascisti chiamati republichini che sostengono la rep. di Salo) \\
Ma è stata letta anche come guerra di classe \t guerra contro i ceti alti e il mondo capitalista che aveva sostenuto mussolini \t la resistenza rossa legge questa lotta anche in chiave rivoluzionaria

\v

Nel 43 si creano le prime bande partigiane, in modo clandestino \t e le fila della resistenza diventano numerosi \t in particolare si uniscono giovani che avevano rifiutato la leva militare \t si arriva a 10.000 militari \\
Questi gruppi si aggregheranno anche in base politica:
\bi
\item brigate garibaldi: + numerosi e comuniste
\item birgate matteoti erano socialisti
\item brigate del popolo erano cattoliche
\item brigate della tradizione repubblicana?
\ei
Queste sono formazioni politiche, poi pero ci sono gruppi a se statni che avevano unico obbettivo di combattere fascimo senza connotazione politica

\v

Poi nasce a Roma il CLN (comitato di liberazione nazionale) = comitato partigiano che deve cordinare l azione delle varie bande \t quello di roma è nazionale \\
Al nord ne crescono diversi ma locali, e il principale sara quello di milano \t il CLN di milano diventa CLNAI (CLN dell'alta italia) \\
Al CLN aderiscono partito 
\bi
    \item comunista
    \item socialista
    \item liberale
    \item democrazia cristiana
\ei
Sono i grandi partiti dle dopo guerra, che saranno antifascisti \\
Il CLN era partigiano, ma diventa formaizone politica con unione dei partiti \t e unione è garantita dalla lotta antifascista \\
C'è pero grande eterogeneita all'interno \t i comunisti volevano una rivoluzione sovietica dopo il fascismo, mentre i socialisti volevano politica riformista \t mentre democristiani volevano stato democratico che permetta ai cattolici partecipazione politica \\
Combattono il fascimo ma non in contesto di comunanza e alleanza piena

\v

Il CLN si confronta prima di tutto su quesitone istituzionale \t i comunisti e socialisti vogliono subito repubblica, perche monarchia si è macchiata di connivenza con mussolini \\
Liberali e cattolici vogliono ora invece la monarchia, e questo è cio che volevano anche gli americani \\
La situzione si sblocca grazie a Togliatti (segretario del parito comunista) \t che tiene discorso a Salerno = la "Svolta di salerno", che è capitale provvisoria \t qua dice che bisogna unire le forze per liberare dai fascisti \\
A guerra finita si fara referndum, pero a guerra finita \t ora non si deve prendere decisione ma liberare l'italia \\
Forze condividono questo compromesso \t quindi Emanuele III affida i poteri a Umberto provvisoriamente, che in realta non diventa re ma ???? \t perche il re sa di essere colluso con il fascismo, e spera di salvare la monarchia affidando il potere al figlio \\
C'è tentativo estremo di salvere la monarchia \t questo nel 44 \\
Nel giugno 44 roma viene liberata dagli americani \t a questo punto Badoglio si dimette, e viene formato nuovo governo con a capo Bonomi (socialista, che aveva fondato Democrazia del lavoro partito) \\
Lui lavora con il CLN \t obbiettivi:
\bi
\item eliminare fascismo (defascistizzazione) dello stato 
\item aiutare resistenza al nord
\item rendersi autonomi dall'amministrazione alleata 
\ei
Infatti gli alleati stavano amministrando le terre che liberavano \t l'amministrazione alleata sara presente finche situazione non sara controllata, e si vuole pero limitare questa azione americana \\
Programma non viene attuato per intero \t anche perche c'è frammentazione delle diverse forze interne, e gli alleati non vogliono cedere

\v

Nel frattempo Repubblica Salo è in crisi \t tutte le forze economiche prendono contatti con gli alleati (anche quelle che avevano fatti accordi con fascismo) \\
Il capitalismo italiano va da chi sta vinceno, quindi con alleati \\
Ed era chiaro che repub. salo fosse fantoccio \t mussolini non conta piu nulla infatti \\
La repubblica di salo vive solo con appoggio nazista, e nazisti sono presenti a tutti i livelli dell amministrazione della repubblica \\
Quindi si inasprisce anche la persecuzione contro ebrei \t nel novembre 43 era stato ordinato che tutti gli ebrei dovevano essere internati a Ponte ??? \\
Con la caduta del fascismo si apre periodo + buio della persecuizone ebrei in italia \t si salvano solo al sud \\
Al centro nord invece ci sono rastrellamenti ovunque, in paesi e citta \t nell ottobre 43 vengono catturati mille ebrei mentre era stato firmato l armistizio, ma roma non era stata liberata \\
Repubblica salo mette a punto anche nuovo progrmma antisemita \t che afferma che ebrei sono tutti stranieri, e quindi in guerra sono nemici \\
Campo di San sabba  e Fossoli (aperto nel 42 per priogionieri di guerra inglesi) ma nel 43 diventa campo di concentramento, e nel 44 passa sotto diretto controllo dell' SS \t fossoli si trovava lungo linea ferroviaria per la polonia \\
Da fossoli vengono deportate 5000 persone, di cui solo la meta ebrei \t tutti gli altri erano oppositori \\
La risiera di San Sabba (era una fabbrica di riso) e diventa campo di sterminio con camere a gas e tutto \\
Ci sono anche campi a Cuneo e a Bolzano

\v

Nel 44 il numero di partigiani aumenta e l azione diventa importante \t nelle vallate dell appennino conquistano posizioni importanti \\
Nasce la repubblica di Montefiorino, che tiene testa all'offensia tedesca per un mese \t prima repubblica partigiana \\
In val dossola anche, e addirittura qua riaprono scuole e sindacati \\
Nel 44 viene liberata firenze solo grazie a partigiani \t in cui pero le rappresagli tedesche sono fuoriose \\
A roma vengono uccisi 30 soldati tedescchi \t in risposta con le fosse Andreatine fucilano 300 perosne \\
Poi strage di Marzabotto \t a Bologna 1000 civili vengono uccisi dai tedeschi \\
Nel 44 resistenza viene unificata sotto la guida di Cadorna (governo di roma), Parri (azione) e Longo (comunista) \t e Parri guidera il primo governo del dopoguerra, che durera poco pero \\
Unificazione del movimento partigiano da vita al CVL (corpo volontari della liberta) = riconosciuto sia da gov. di roma sia da alleati e diventa forza politica e militare riconosciuta \\
Poi crisi, perche tra il 44 e 45 la resistenza vive momento difficile \t alleati stavano risalendo e vengono fermati sulla linea gotica per tutto l'inverno del 44 \\
Fermi qui, gli alleati e comandante Alexander a novemrbe sospendono offensiva contro nazzisti sull appennino toscano \t il fronte di combattimento principale diventa la linea gotica, e quindi il fronte sull appennino deve essere sospeso \\
Questo significa che partigiani non vengono + supportati \t e partigiani non ubbidiscono e continuano lotta per liberazione = scelta estrema \\
GLi aiuti alleati non sono sufficienti, e i tedeschi attuano una controoeffensiva importante che travole i luoghi liberati \t sanno che partigiani non sono supportati \\
Alla primavera del 45 alleati riprendono offensiva, sfondano la linea e riforniscono partigiani \t e intanto il numero dei partigiani aumenta \t i tedeschi sono in ritirata \\
Ci si prepara alla conclusione della guerra \\
Il 21 di aprile alleati entrano a bologna \t genova e milano si liberano da sole, e il 25 aprile Milano è libera \\
Entro la fine di aprile viene liberata anche Torino

\v

Mussolini prova a scappare \t si traveste da soldato tedesco, ma viene riconosciuto e portato a Dongo insieme alla Petacci e viene fucilato \\
Poi viene impiccato per i piedi ed esposto in piazzale loreto \\
Il 30 aprile Hitler si suicida nel bunker di berlino \\
In alcune zone del nord italia si continua a combattere per qualche gionro \t e poi italia è libera \\
Non finisce pero la forza armata \t c'è problema della riconsegna delle armi \t infatti per mesi ancora i partigiani uccidono fascisti, iscritti al partito fascista e militanti \\
Succede che ci si vendica \t viene ordinata la riconsegna delle armi, quindi la gente si tiene le armi e si vendicano \\
Ci sono state delle esecuzioni anche barbare, e poi venivano saccheggiate la casa e i beni \t chiamata fase della giustizia inurrezionale

\s[Ricostruzione]
Si vuole ora ricostruire, materialmente e anche moralmente \\
Si devono ricostuire case, industrie, ponti, scuole etc. \t la produzione industriale era un terzo \\
Beni di sussitenza mancano, case mancano e costo della vita alto \t anche disoccupazione alta \\
Nel 1946 l'amministrazione delle nazioni unite per la ricostruzione si impegna a fornire all italia macchinari, cibo e materie prime all italia \t ma aiuti non sono sufficienti, quindi piano Marshall \\
Il piano Marshall è piano di ricostruzione proposto dagli stati uniti \t europa deve tornare in piedi, perche è partner commerciale importante \\
C'è anche un pensiero ideologico pero \t nel 47 c'è Truman come presidente, che vuole combattere comunismo a tutti i livelli \t quindi Marshall (ministro interni???) propone piano per far vedere che paesi occidentali sono migliori di quelli sovietici \\
Il piano marshall viene in realta proposto anche a stati dell est e URSS, che ovviamente rifiutano \\
Il piano marsh. sara anche la condizione per la nascita del Patto Atlantico \t usa aiuta europa, che quindi è un po indebitata \\
USA stanzia 14 miliardi \t enra in vigore nel 51?? \\
Il piano marshall viene considerato da alcuni come definitivo, secondo altri è stato inutile ma è servito al governo per contenere spesa pubblica e ottenere bilancio \t non è servito per il rilancio del paese

\v

Fino al maggio del 47 i governi italiani includono tutti \t Governi di Unita antifascita \t sono tutte le forze al governo, che quindi non sono coesi \t si vuole pero abbandonare tutto cio che è fascista, volonta comune \\
Tutti vogliono mercato libero \t abbandonare fascismo economico \\
Einaudi diventa presidente del consiglio nel 47 ed era econmista liberale \t svalutera la lira e punta sulla stabilizzaizone della lira e abbassare inflazione \t diventera anche pres. repubblica \\
Dal 47 cambia \\
Nel 45 si ha primo govenro presieduto da Ferruccio Parri \t capo della resistenza, ma dura poco \t parri era di formazione socialista \\
Cosi si delineano due schieramenti: da una parte democrazia cristiana con USA e ceto medio, dall latro sindacati, CGL, proletariato, URSS \\
Contrapposizione pero non degnera \t Togliatti dice che non è il momento di fare la rivoluzione, ma non è momento \t non c'è neanche trattato di pace \\
Togliatti insieme a Nelli chiedono elezioni, che ancora non ci sono state \t popolo si deve esprimere \t e elezioni anche per assemblea costituente \\
De Gasperi (leader DC) rimanda le elezioni al 46, e mette d accordo le forze politiche varie \t lui da sua impronta al parito \\
DC diventa partito democratico moderno e interclassista \t si rivolge a tutte le classi, ed è partito confessionale \\
Il mondo imprenditoriale supporta dc \t perche c'è paura delle sinistre con rischio rivoluzione \t DC supportata anche dalla chiesa \\
De gasperi ha programma quindi di chiesa, libero mercato, anticomunismo \t punti fond. \t governera fino al 53, e il primo governo de gasperi vede togliatti ministro di grazia e giustizia, e nenni di ? \\
Tutte le forze diverse sono ancora la governo

\v

De Gasperi genera proprio un contesto per la democrazia cristiana che le permettera di governare \\
Vuole prima di tutto al voto \t referendum popolare per monarchia o repubblica \t 2 giugno del 46 c'è referendim, votano donne (prime elezioni a suff. univ.) e vince la repubblica \\
Rep. vince al nord e al centro, mentre al sud vince monarchia \t con scarto minimo, quindi accusato broglio elettorale \\
Umberto II va quindi in esilio di Portogallo, e il padre nel frattempo aveva abdicato in favore di Umb. II \\
Umb. II viene chiamato il re di maggio \t perche è stato re solo a maggio, padre aveva abdicato a inizio maggio e fine a giugno \\
L'elezioni per l assemblea costituente si tengono in parallelo, e vincono partiti di massa: dc e altri due
\end{document}
